%! Vectors — Unit 4, Lesson 4.8 — Lesson Plan
\documentclass[10pt]{article}
\usepackage{precalculus-article}
\usepackage{precalculus-boxes}
\usepackage{graphicx}
\graphicspath{{images/}}
\newcommand{\TallMath}[1]{$\displaystyle #1\rule[-1.4em]{0pt}{3.2em}$}
\newcommand{\CourseName}{AP Precalculus}
\newcommand{\SchoolYear}{2026--2027}
\newcommand{\MeetingLength}{55 minutes}
\newcommand{\UnitNumberName}{Unit 4: Functions Involving Parameters, Vectors, and Matrices \quad}
\newcommand{\LessonNumberName}{Lesson 4.8: Vectors}

\begin{document}

%----------------------------------------------------------------------
% Title Block
%----------------------------------------------------------------------
\begin{center}
  {\Large\bfseries \CourseName: \SchoolYear} \\[4pt]
  {\small\bfseries \UnitNumberName \LessonNumberName \quad (3 instructional periods, \MeetingLength\ each)}
\end{center}

%----------------------------------------------------------------------
% Primary Objective
%----------------------------------------------------------------------
\begin{tcolorbox}[colback=sky, colframe=navy, boxrule=0.9pt, arc=2mm,
    left=2mm, right=2mm, top=1.5mm, bottom=1.5mm]
  \textbf{Primary Objective:} Students will identify characteristics of vectors
  (magnitude, direction, component form), perform scalar multiplication, vector
  addition, and dot products, find unit vectors, and use the dot product formula
  to determine the angle between two vectors and test for perpendicularity.
  \textit{(Big Idea: Functions)}
\end{tcolorbox}

\vspace{0.08in}

%----------------------------------------------------------------------
% Priority Ideas & Skills
%----------------------------------------------------------------------
\begin{skillbox}[Priority Ideas \& Skills]{goldbox}
\begin{tabularx}{\linewidth}{@{} p{0.46\linewidth} X @{}}
\textbf{AP Skills} & \textbf{Key Understandings (EKs)} \\
\midrule
\textbf{2.A} \textit{Identify information from representations} ---
  Read the tail, head, component form $\langle a,b\rangle$, magnitude, and
  direction of a geometrically represented vector; use trigonometry to
  extract components when magnitude and direction angle are given.
&
\textbf{EK 4.8.A.1--A.4}: A vector is a directed segment with tail $P_1$
and head $P_2$. Components $\langle a,b\rangle = \langle x_2-x_1,y_2-y_1\rangle$.
Magnitude $= \sqrt{a^2+b^2}$. Direction parallel to segment from origin to
$(a,b)$; components follow from $a = r\cos\theta$, $b = r\sin\theta$. \\[4pt]

\textbf{3.A} \textit{Describe characteristics} ---
  Explain the geometric effect of scalar multiplication (parallel, scaled
  magnitude) and vector addition (head-to-tail); describe the dot product as
  a scalar encoding the angle between two vectors.
&
\textbf{EK 4.8.B.1--B.3}: Scalar multiplication: $k\langle a,b\rangle
= \langle ka,kb\rangle$. Addition: add corresponding components; graphically,
head-to-tail. Dot product: $\langle a_1,b_1\rangle\cdot\langle a_2,b_2\rangle
= a_1 a_2+b_1 b_2$ (a scalar). \\[4pt]

\textbf{3.B} \textit{Apply numerical results} ---
  Compute unit vectors; express any vector in terms of $\hat{\imath}$ and
  $\hat{\jmath}$; apply the dot product formula to find the angle between two
  vectors; identify perpendicular vectors by testing whether the dot product
  equals zero.
&
\textbf{EK 4.8.C.1--C.2, 4.8.D.1--D.2}: Unit vector:
$\hat{u} = \vec{v}/\|\vec{v}\|$. Standard basis:
$\langle a,b\rangle = a\hat{\imath}+b\hat{\jmath}$.
Dot product $= \|\vec{v}_1\|\,\|\vec{v}_2\|\cos\theta$; perpendicular iff
dot product $= 0$. Law of Cosines/Sines applies to vector-addition triangles. \\
\end{tabularx}
\end{skillbox}

\vspace{0.08in}

%----------------------------------------------------------------------
% Vocabulary
%----------------------------------------------------------------------
\begin{skillbox}[{Vocabulary, Concepts, \& Theorems}]{greenbox}
\renewcommand{\arraystretch}{1.4}
\begin{tabularx}{\linewidth}{@{} p{0.24\linewidth} X @{}}
\textbf{Term} & \textbf{Definition / Note} \\
\midrule
vector &
  A directed line segment described by magnitude and direction; written
  $\vec{v} = \langle a, b\rangle$ in component form. \\[4pt]
tail / head &
  Starting point (tail) and ending point (head) of the directed segment. \\[4pt]
magnitude &
  \TallMath{\|\vec{v}\| = \sqrt{a^2+b^2};} the length of the vector. \\[4pt]
scalar multiplication &
  $k\langle a,b\rangle = \langle ka,kb\rangle$; scales the magnitude by $|k|$,
  reverses direction when $k<0$. \\[4pt]
dot product &
  \TallMath{\langle a_1,b_1\rangle \cdot \langle a_2,b_2\rangle
  = a_1 a_2+b_1 b_2;} a scalar, not a vector. \\[4pt]
unit vector &
  Magnitude 1; $\hat{u} = \vec{v}/\|\vec{v}\|$.
  Standard basis: $\hat{\imath}=\langle 1,0\rangle$,
  $\hat{\jmath}=\langle 0,1\rangle$. \\
\end{tabularx}
\end{skillbox}

\vspace{0.08in}

%----------------------------------------------------------------------
% Activate Prior Knowledge
%----------------------------------------------------------------------
\begin{skillbox}[Activate Prior Knowledge \& Spiral Review]{sky}
\begin{tabularx}{\linewidth}{@{}X c@{}}
\begin{minipage}[t]{0.96\linewidth}\vspace{0pt}
\textbf{Spiral topics activated during Warm-Up (first 8 minutes):}
\begin{itemize}[leftmargin=1.4em, itemsep=2pt]
  \item \textbf{Distance formula} (Algebra 2): $\sqrt{(x_2-x_1)^2+(y_2-y_1)^2}$ ---
    this is exactly the magnitude formula once components are named.
  \item \textbf{Trigonometry on the unit circle} (Lesson 4.4):
    $\cos\theta$ and $\sin\theta$ as coordinates; needed to extract vector
    components from a magnitude-direction description.
  \item \textbf{Parametric component functions} (Lesson 4.2):
    evaluating pairs of expressions at specific values --- warmup for
    computing vector expressions under scalar multiplication.
\end{itemize}
\end{minipage}
&
\end{tabularx}
\end{skillbox}

\vspace{0.08in}

%----------------------------------------------------------------------
% Hook
%----------------------------------------------------------------------
\begin{skillbox}[Hook --- Entry Event]{sky}
\textbf{Scenario:} A drone leaves point $A=(1,2)$ and arrives at point $B=(5,6)$.
A second drone leaves $C=(3,0)$ and arrives at $D=(7,4)$. Both drones are programmed
with \emph{identical} displacement instructions.

\textbf{Discussion prompts:}
\begin{enumerate}[leftmargin=1.4em, itemsep=3pt]
  \item \textit{``How far did each drone move east? How far north? What does
    it mean that both drones received the same instructions even though they
    started at different points?''}
  \item \textit{``If the engineer doubles the flight instructions for drone A,
    where does it end up? What happens to the distance traveled?''}
\end{enumerate}

\textbf{Key tension:} Position is described by coordinates; \emph{displacement}
is described by direction and magnitude --- a vector. Vectors encode relative
change, not absolute location. That distinction drives the entire lesson.
\end{skillbox}

\vspace{0.08in}

%----------------------------------------------------------------------
% Lesson
%----------------------------------------------------------------------
\begin{skillbox}[Lesson --- Instruction Sequence (Periods 1--3)]{sky}
\begin{multicols}{2}

\textbf{Step 1 --- Vector Notation and Components} (12 min, Period 1)

Define: vector = directed segment from tail $P_1=(x_1,y_1)$ to head
$P_2=(x_2,y_2)$. Components $a=x_2-x_1$, $b=y_2-y_1$; write
$\vec{v}=\langle a,b\rangle$. Zero vector $\langle 0,0\rangle$ when
$P_1=P_2$. Magnitude: $\|\vec{v}\|=\sqrt{a^2+b^2}$.

Direction from trig: magnitude $r$, angle $\theta$ $\Rightarrow$
$a=r\cos\theta$, $b=r\sin\theta$.

Worked example: $P_1=(1,3)$, $P_2=(5,6)$:
$\vec{v}=\langle 4,3\rangle$, $\|\vec{v}\|=5$. (EK 4.8.A.1--A.4)

\columnbreak

\textbf{Step 2 --- Scalar Multiplication \& Vector Addition} (13 min, Period 1)

Scalar: $k\langle a,b\rangle=\langle ka,kb\rangle$. New magnitude: $|k|\|\vec{v}\|$.
Parallel to original; reversed when $k<0$.

Addition: $\langle a_1,b_1\rangle+\langle a_2,b_2\rangle
=\langle a_1+a_2,b_1+b_2\rangle$.
Geometry: head-to-tail; resultant from first tail to final head.

Example: $\vec{u}=\langle 3,1\rangle$, $\vec{w}=\langle -1,4\rangle$.
$\vec{u}+\vec{w}=\langle 2,5\rangle$; $2\vec{u}=\langle 6,2\rangle$.
(EK 4.8.B.1--B.2)

\end{multicols}

\begin{multicols}{2}

\textbf{Step 3 --- Dot Product and Unit Vectors} (15 min, Period 2)

Dot product: $\langle a_1,b_1\rangle\cdot\langle a_2,b_2\rangle
=a_1 a_2+b_1 b_2$ --- a scalar.

Unit vector: $\hat{u}=\dfrac{\vec{v}}{\|\vec{v}\|}
=\left\langle\dfrac{a}{\|\vec{v}\|},\dfrac{b}{\|\vec{v}\|}\right\rangle$.

Standard basis: $\hat{\imath}=\langle 1,0\rangle$, $\hat{\jmath}=\langle 0,1\rangle$;
every vector $\langle a,b\rangle=a\hat{\imath}+b\hat{\jmath}$.

Example: $\vec{v}=\langle 3,4\rangle$, $\|\vec{v}\|=5$,
$\hat{u}=\langle 3/5,4/5\rangle$. (EK 4.8.B.3, 4.8.C.1--C.2)

\columnbreak

\textbf{Step 4 --- Angle Between Vectors \& Perpendicularity} (15 min, Period 2)

Formula: $\vec{v}_1\cdot\vec{v}_2=\|\vec{v}_1\|\,\|\vec{v}_2\|\cos\theta$.

Solve: $\theta=\cos^{-1}\!\left(\dfrac{\vec{v}_1\cdot\vec{v}_2}
{\|\vec{v}_1\|\,\|\vec{v}_2\|}\right)$.

Perpendicular $\Leftrightarrow$ dot product $=0$ (since $\cos 90^\circ=0$).

Example: $\vec{u}=\langle 1,2\rangle$, $\vec{w}=\langle 4,-2\rangle$.
$\vec{u}\cdot\vec{w}=4-4=0$ $\Rightarrow$ perpendicular.

Law of Cosines connects $\|\vec{u}+\vec{w}\|$ to the triangle with sides
$\|\vec{u}\|$, $\|\vec{w}\|$, and included angle. (EK 4.8.D.1--D.2)

\end{multicols}

\smallskip
\textbf{Pacing note:} Steps 1--2 fill Period 1; Steps 3--4 fill Period 2;
Period 3 is activity, exit ticket, and homework launch.
\end{skillbox}

\vspace{0.08in}

%----------------------------------------------------------------------
% Active Monitoring
%----------------------------------------------------------------------
\begin{skillbox}[Active Monitoring]{redbox}
\begin{itemize}[leftmargin=1.4em, itemsep=3pt]
  \item \textbf{Confusing position with displacement}: Students write the
    head coordinates as the vector. Prompt: \textit{``A vector is the
    \emph{change} in position --- subtract tail coordinates from head coordinates.''}
  \item \textbf{Dot product treated as a vector}: Students try to write
    $\vec{u}\cdot\vec{w}=\langle\ldots\rangle$. Reinforce: the dot product
    is a \emph{number}.
  \item \textbf{Unit vector partial division}: Students divide only one
    component by the magnitude. Prompt: \textit{``Every component is divided
    by $\|\vec{v}\|$ --- what is the magnitude of your answer? It must equal 1.''}
  \item \textbf{Perpendicularity and zero dot product}: Students believe only
    zero vectors can have a zero dot product. Ask: \textit{``Can two nonzero
    vectors be perpendicular? Give me a concrete example using
    $\hat{\imath}$ and $\hat{\jmath}$.''}
\end{itemize}
\end{skillbox}

\vspace{0.08in}

%----------------------------------------------------------------------
% Group Work & Differentiation
%----------------------------------------------------------------------
\begin{skillbox}[Group Work \& Differentiation]{redbox}
Students work in groups of 3--4 on the Navigation Mission activity (tiered).

\begin{itemize}[leftmargin=1.4em, itemsep=4pt]
  \item \textbf{Tier R (Remediate)}: Given $\vec{u}=\langle 3,4\rangle$ and
    $\vec{w}=\langle -3,4\rangle$, compute $\|\vec{u}\|$, $\|\vec{w}\|$,
    $\vec{u}+\vec{w}$, and $\vec{u}\cdot\vec{w}$. State whether $\vec{u}$ and
    $\vec{w}$ are perpendicular.
  \item \textbf{Tier A (Approaching Proficiency)}: A ship travels displacement
    $\vec{u}=\langle 6,8\rangle$ km, then $\vec{w}=\langle -2,4\rangle$ km.
    Find the total displacement vector, its magnitude, and the unit vector
    in the direction of the total displacement.
  \item \textbf{Tier E (Extension)}: Find all vectors $\langle a,b\rangle$ of
    magnitude 5 that are perpendicular to $\langle 3,4\rangle$. Use the dot
    product condition and magnitude condition simultaneously. How many solutions
    exist? Generalize to any nonzero vector.
\end{itemize}
\end{skillbox}

\vspace{0.08in}

%----------------------------------------------------------------------
% Individual Work & Assessment
%----------------------------------------------------------------------
\begin{skillbox}[Individual Work \& Assessment]{redbox}
Exit ticket administered independently (final 8 minutes of Period 3).

\begin{enumerate}[leftmargin=1.4em, itemsep=4pt]
  \item Find the component form and magnitude of the vector from $P_1=(2,5)$
    to $P_2=(6,2)$.
  \item Given $\vec{u}=\langle 4,3\rangle$ and $\vec{w}=\langle 1,-2\rangle$,
    compute $\vec{u}\cdot\vec{w}$. Are the vectors perpendicular? Explain.
  \item Find the unit vector in the direction of $\vec{v}=\langle 5,12\rangle$.
\end{enumerate}

\textbf{Data use:} Students who cannot find components without prompting need
additional practice distinguishing position from displacement; revisit
EK 4.8.A.2 at the start of the next lesson. Students who compute the dot
product but cannot connect it to perpendicularity need a worked example
linking $\cos 90^\circ=0$ to the angle formula.
\end{skillbox}

\vspace{0.08in}

%----------------------------------------------------------------------
% Reinforcement & Extension
%----------------------------------------------------------------------
\begin{skillbox}[Reinforcement \& Extension]{goldbox}
\textbf{Homework (Lesson 4.8):} 5--6 problems covering: finding components
and magnitude from given endpoints; scalar multiplication and vector addition;
unit vectors; dot product; perpendicularity; one problem applying the Law of
Cosines to find the magnitude of a vector sum from two given vectors and the
angle between them.

\textbf{Extension:} For two vectors $\vec{u}$ and $\vec{w}$, what does
$\vec{u}\cdot\vec{w}=\|\vec{u}\|\,\|\vec{w}\|$ imply geometrically? What does
$\vec{u}\cdot\vec{w}=-\|\vec{u}\|\,\|\vec{w}\|$ imply? At what angle do they
point relative to each other in each case?

\textbf{Preview of Lesson 4.9:} Ask students:
\textit{``You now know how to add vectors. What if we organized multiple
vectors as rows in a rectangular array? Could we define an operation
analogous to multiplication for such arrays?''} This leads to matrices.
\end{skillbox}

\end{document}
